\documentclass[10pt,a4paper]{amsart}
\usepackage[margin=19mm]{geometry}
\usepackage{amsmath,amssymb,mathtools}
\usepackage[scr=rsfs,cal=euler]{mathalfa}
\usepackage{xfrac}
\usepackage[unicode,hidelinks,bookmarks=false]{hyperref}
\newcommand{\ie}{i.e.\ }
\newcommand{\cf}{cf.\ }
\newcommand{\resp}{respectively\ }
\newcommand{\Subsection}{\subsection}
\providecommand{\nobreakdash}{\nobreak\hbox{-}}
\setlength{\emergencystretch}{2em}
\multlinegap0pt
\usepackage{bm}
\usepackage{bbm}
\usepackage{dsfont}
\usepackage{xspace}
\usepackage{cancel}
\usepackage{mathtools}
\usepackage{amsthm}
\makeatletter
\@ifundefined{theorem}{\newtheorem{theorem}{Theorem}}{}
\@ifundefined{proposition}{\newtheorem{proposition}[theorem]{Proposition}}{}
\makeatother
\makeatletter
\expandafter\ifx\csname construction\endcsname\relax
\newenvironment{construction}    % this is the environment name for the input
{\renewcommand{\qedsymbol}{$\lozenge$}%
	\pushQED{\qed}\begin{constr}}
	{\popQED\end{constr}}
\fi
\expandafter\ifx\csname customthm\endcsname\relax
\newenvironment{customthm}[1]
{\renewcommand\theinnercustomthm{#1}\innercustomthm}
{\endinnercustomthm}
\fi
\makeatother
\title[Logarithmic Quot spaces, boundedness, and K-tropicalizations]{Logarithmic Quot spaces, boundedness, and K-tropicalizations}
\author{Patrick Kennedy-Hunt, Dhruv Ranganathan}
\date{Source: arXiv:2508.08175v1 (CC BY 4.0)}
\begin{document}
\maketitle

\noindent\emph{Source and licence.} Logarithmic Quot spaces, boundedness, and K-tropicalizations. Version arXiv:2508.08175v1 (CC BY 4.0), distributed under the Creative Commons Attribution 4.0 International licence, as recorded in the arXiv licence metadata for this exact version and shown on the abstract page https://arxiv.org/abs/2508.08175v1. Full rights notice: licenses/arxiv-2508.08175-CC-BY-4.0.txt.\par\smallskip

\noindent\textit{Editorial scope.} Counted unit: the Proposition whose source label is \texttt{prop:Toric Case Enough} in the article (source0.tex lines 2042--2044, source label \texttt{prop:Toric Case Enough}), delivered together with the article's own complete original proof of it (source0.tex lines 2046--2073, one \texttt{proof} environment of 28 lines). No mathematics is written, repaired or re-derived: statement and proof are the article's own text line for line. Required context retained uncounted: thm: combinatorial-boundedness (lines 2020--2022); prop: global-toric-embedding (lines 2026--2032). Every definition, display and label of those lines is retained unchanged. Omitted from this package and not redistributed: 1--2019, 2023--2025, 2033--2041, 2045--2045, 2074--2339. Nothing inside a retained excerpt is omitted or altered: every retained body is delivered exactly as the article's lines are written, so no omission receipt of any kind accompanies this package. Citations: the retained excerpts carry no citation command at all, so no bibliography entry of the article is transcribed and no reference is redistributed. Reference adaptation: none was needed and none was made; every reference inside the delivered unit resolves inside it, so no unresolved reference or citation is produced. Printed slips kept literal: no printed word, symbol, hypothesis or number of the counted statement or of its proof was altered. Layout and class: the article's own class is \texttt{amsart} with packages including subfiles, lipsum, quoting, tabu, fullpage, diagbox, booktabs, placeins, geometry, xcolor, graphicx, amssymb, csquotes, dsfont; the wrapper is the lane's pinned \texttt{amsart} editorial wrapper at 10pt on A4, which restates the theorem-like environments the retained text needs and adds the article's own macro definitions that the retained text needs (none), with extra wrapper packages (bm, bbm, dsfont, xspace, cancel, mathtools). The delivered numbering is the unit's own. Assets: no retained span contains a figure environment, a graphics-inclusion command, a PDF-page inclusion, a table or a TikZ picture; no image file is copied into this package, the package declares no asset dependency and no image file is redistributed with it. Licence: version arXiv:2508.08175v1 of this article is distributed under the Creative Commons Attribution 4.0 International licence, as recorded in the arXiv licence metadata for this exact version and shown on the abstract page https://arxiv.org/abs/2508.08175v1. A full rights notice is delivered at licenses/arxiv-2508.08175-CC-BY-4.0.txt. Characters: every retained excerpt is pure ASCII, so the delivered engine, XeLaTeX with its default Latin Modern text font, reports no missing character.
\begin{theorem}\label{thm: combinatorial-boundedness}
The set of combinatorial types of $K$-tropicalizations of points of $\mathsf{Quot}_\Lambda(X|D,\mathcal E)$ is finite. 
\end{theorem}

\begin{proposition}\label{prop: global-toric-embedding}
There is a strict embedding
\[
X\hookrightarrow Y
\]
where $Y$ is the projectivization of a direct sum of line bundles over a product of projective spaces $\mathbb P^{\underline n}$, where $Y$ is equipped with its fiberwise toric logarithmic structure, denoted $\partial Y$.
\end{proposition}

\begin{proposition}\label{prop:Toric Case Enough}
Assume Theorem~\ref{thm: combinatorial-boundedness} for pairs $(Y,E)$ consisting of a toric variety and its toric boundary. Then Theorem~\ref{thm: combinatorial-boundedness} holds for all pairs $(X,D)$. 
\end{proposition}

\begin{proof}
Consider tropicalizations of points in the moduli space $\mathsf{Quot}_\Lambda(X|D,\mathcal E)$. We exhibit an injective map from the set of tropicalizations arising from $\mathsf{Quot}_\Lambda(X|D,\mathcal E)$ into a corresponding set of tropicalizations arising from a toric pair, with controlled asymptotics $K$-weights. 

We start with the geometric input from Proposition~\ref{prop: global-toric-embedding}. Let $Y$ be the toric variety arising in the proposition, and 
\[
X\hookrightarrow Y
\]
the strict embedding. Recall that $Y\to \mathbb P^{\underline n}$ is a projective space bundle and we have equipped $Y$ with a fiberwise toric logarithmic structure. 

Consider a point of $\mathsf{Quot}_\Lambda(X|D,\mathcal E)$. It is represented by an algebraically transverse sheaf on an expansion $X_\Gamma$ of $X$ along $D$. Since $X\hookrightarrow Y$ is strict we can express $X_\Gamma$ as a pullback of an expansion
\[
Y_\Gamma\to Y
\]
along $j\colon X\hookrightarrow Y$. A logarithmic quotient $\mathcal F$ of $\mathcal E$ on $X_\Gamma$ pushes forward to a logarithmic quotient $j_\star\mathcal F$ of $j_\star\mathcal E$ on $Y_\Gamma$. Consider the ample line bundle $\mathcal O(1,\ldots,1)$ on $\mathbb P^{\underline n}$. Using the projection formula in $K$-theory, the $K$-weights of this quotient with respect to this ample and all the line bundles making up the constituent line bundles of $Y\to \mathbb P^{\underline n}$ is determined by the $K$-weights of $\mathcal F$.  

We now enlarge the logarithmic structure on $Y$ from $\partial Y$ to make it toric. Fix an algebraically sheaf $\mathcal F$ on $Y_\Gamma\to Y$. Consider the base $\mathbb P^{\underline n}$ of the toric bundle. On a $\mathbb P^{n_i}$ factor of the product, choose $(n_i+1)$ generic hyperplanes. Each choice of these hyperplanes makes $\mathbb P^{n_i}$ isomorphic to a toric variety, with a different torus for each choice. 

Similarly, choosing such a set of generic hyperplanes in each factor, we turn $\mathbb P^{\underline n}$ into a toric variety, and by pulling back $Y\to \mathbb P^{\underline n}$, such a choice of hyperplanes makes $Y$ into a toric variety as well. We refer to this system of hyperplane sections as a {\it system of residual toric divisors}. Note that different choices of toric divisors give rise to isomorphic toric varieties, in fact, they are related by a factorwise projective change of coordinates. 

\noindent
{\bf Claim.} Let $\mathcal H$ be a generic system of residual toric divisors. The sheaf $\mathcal F$ on $Y_\Gamma$ is algebraically transverse with respect to the union of $\partial Y$ and $\mathcal H$. 

The claim follows from a Bertini argument -- given $\mathcal F$ algebraically transverse to $\partial Y$, a generic section of a basepoint free line bundle on $Y$ not contain any sections of $\mathcal F$, and the same is true on any stratum. Repeating this for all the elements forming a system of residual toric divisors proves the claim. We can do this for every for every logarithmic quotient. 

Putting this together, we obtain, for each quotient $\mathcal F$ on $X_\Gamma$, a quotient on an expansion $Y_\Gamma$ of a toric variety, along its full toric boundary, with predictable asymptotic $K$-weights depending only on the starting from the asymptotic $K$-weights of $\mathcal F$. The tropicalization of the sheaf with respect to this larger logarithmic structure $E$ certainly determines the one respect $\partial Y$. We emphasize here that that the choice of residual divisor system $\mathcal H$ will depend on $\mathcal F$, by changing coordinates as above, we can view each of the resulting quotient as an object on a fixed toric variety. 

We have therefore constructed an injective map from the set of all tropicalizations arising from the $(X,D)$-problem into the set of all tropicalizations arising in the $(Y,E)$ problem, and the latter is toric, so the proof of the proposition is complete. 
\end{proof}

\end{document}
